\scanpage{85}
ce qui implique en particulier que $B_T\in\mathrm{Ob}\,\Delta$ (si $\Delta$ est défini par des propriétés d'exactitude, cela est clair, car dans $B_T=\underline{\mathrm{Hom}}(T,\tau)$ les $\underset{\unsure{\leftrightarrow|\leftarrow}}{\lim}$ de type $d',d''$ se construisent argument par argument)\unsure{;|}. Mais il y a un ``grain de sel'', c'est qu'il n'est pas clair que $B_T$ soit une $\mathcal U$-catégorie.

La description de $B_T$ qu'on a donné peut sembler brutale ; on va définir pas à pas une catégorie $B'_T$ de ``constructions explicites'' (sic) pour une \deleted{catégorie}\inserted{théorie} \uline{algébrique} $T$, quand on construit cette notion de théorie algébrique. Mais on supposera que $B'_T\in\mathrm{Ob}\,\Delta$.

\textcircled{a} Structure vide sur $I$ objets : correspond à \deleted{\unsure{$\tau$}} la ``théorie'' $\tau^I$. Elle est représentable par la $\Delta$-enveloppe \deleted{$\widetilde I$} de la catégorie discrète $I$. \deleted{Les objets de $\widetilde I$ sont sans doute obtenus} On pose $B'_{\tau^{I}}=\widetilde I$. On le justifie \inserted{(si $\Delta$ défini par ax. d'exactitude)} en \deleted{\unsure{notant}}\inserted{expliquant} qu'on peut expliciter $\widetilde I$ par lim de types $d$ à partir des objets $i$, et des axiomes d'exactitude en inversant des flèches grâce aux axiomes d'exactitudes

\[
\deleted{\unsure{\Delta\times I\quad I_{\Delta}|I\times I\quad I_{\Delta}}}
\]

\textcircled{b} Structure \unsure{déduit|déduite} d'une structure $T$ \inserted{(\deleted{avec cat. de constr. Explicite}\inserted{avec cat de constructions} $,B'_T$)} en ajoutant des axiomes : \deleted{\unsure{en joignant|en ajoutant}}\inserted{un paquet} de flèches \inserted{$\Sigma$} de $B'_T$ doivent être rendues inversibles\inserted{$B'_{T'}=B'_T\Sigma^{-1}$ (sur $\Delta$)} ;
% margin note written sideways; an arrow connects the note to the axioms.
\margin{il faut aussi prévoir les axiomes du type ``flèches rendues égales'' (en réalité, \unsure{un des|cas des|pas des} précédents si on suppose p.ex. $d'=d''=\varnothing$)}
% slanted
\margin{commentaire sur le type d'axiomes : axiomes d'iso, \inserted{ou égalités de flèches} de \deleted{\unsure{propriété}} \unsure{univers.} d'is., presque \unsure{quelconque} dans \unsure{op.} \unsure{quelques détails p.ex.} \uline{\unsure{sauf le cas exact}}. On peut rajouter des \deleted{\unsure{nouveaux}} objets de base en vrac, mais cel peut être aussi contenu dans \textcircled{d}}
\uline{NB} Si \deleted{on} $T$ est 2-repr. par \deleted{\unsure{$B'$}}$R$, \deleted{donc $B'_T\simeq R$,} alors $T'$ repr. par \deleted{\unsure{$R\Sigma^{-1}$} \unsure{soit}} $B'_{T'}$, \deleted{\unsure{alors} $R'$, car $B'_{T'}\simeq R'$.}

\inserted{Re \uline{NB} Si on a foncteurs fidèles $T\to\tau^I$, le composé $T'\to T\to\tau^I$ aussi, ($T'\to T$ pl. fid)}

\textcircled{c} Structure déduite d'une structure $T$ en ajoutant des flèches (nouvelles données) : \deleted{$\forall\,x,y\in B'_T$, ens $\unsure{H_{xy}|H_{x,y}}$ (\unsure{qui peut être vide})} \unsure{i.e.|On} $H$ un ensemble d'indices, et $H\longrightarrow\mathrm{Ob}\,B'_T\times\mathrm{Ob}\,B'_T$ \quad $h\longmapsto s(h)$, $h\longmapsto b(h)$, et on pose
\[
\begin{aligned}
T'(C)={}&\text{catégorie des objets }x\text{ de }T(C),\\
&\text{\uline{munis} d'une application }\varphi\text{ de}\\
&H\text{ dans }\mathrm{Fl}(C)\text{ telle que }\forall\,h\in H\\
&\varphi(h):s(h)(x)\longrightarrow b(h)(x)
\end{aligned}
\]
% slanted
\margin{\uline{NB} On n'admet pas la donnée de sous-objets, \deleted{\unsure{mais}} on peut avoir \deleted{des} \unsure{nouvel|nouveau} objet de base et un mono de celui-ci !}
% archive sheet number: 84
